\subsection{APPLICATION TO THE BERNOULLI NUMBERS}

Theorem 1 of VI, p.~286, shows that, for \(0 \leqslant x < \pi\) , the series with general term \(\frac{2x}{n^2\pi^2 - x^2} \geqslant 0\) converges. On the other hand, one can write, for every complex number \(z\) such that \(|z| < \pi\) ,

\[
{\frac {2 z}{n ^ {2} \pi^ {2} - z ^ {2}}} = {\frac {2 z}{n ^ {2} \pi^ {2}}} \sum_ {k = 0} ^ {\infty} {\frac {z ^ {2 k}}{n ^ {2 k} \pi^ {2 k}}}
\]

the series on the right being absolutely convergent. We shall deduce from this that the ``double'' series

\[
\sum_ {n = 1} ^ {\infty} \sum_ {k = 1} ^ {\infty} \frac {- 2 z ^ {2 k - 1}}{n ^ {2 k} \pi^ {2 k}}\tag{11}
\]

is absolutely convergent in the open disc \(|z| < \pi\) , normally convergent on every compact set contained in this disc, and has sum \(\cot z - \frac{1}{z}\) . Indeed, for \(|z| \leqslant a < \pi\) the absolute value of the general term of (11) is at most equal to \(2a^{2k-1}/n^{2k}\pi^{2k}\) , and the sum of any finite number of the terms \(2a^{2k-1}/n^{2k}\pi^{2k}\) is less than the finite number \(\sum_{n=1}^{\infty} \frac{2a}{n^2\pi^2 - a^2}\) ; on summing first over \(k\) , then over \(n\) , one sees that the sum of the series (11) is equal to \(\sum_{n=1}^{\infty} \frac{2z}{z^2 - n^2\pi^2}\) , which proves our statement.

If now one sums the series (11) first over \(n\) and then over \(k\) one obtains the identity (for \(|z| < \pi\) )

\[
\cot z - \frac {1}{z} = - 2 \sum_ {k = 1} ^ {\infty} \frac {S _ {2 k}}{\pi^ {2 k}} z ^ {2 k - 1}\tag{12}
\]

where we have put \(S_{k} = \sum_{n=1}^{\infty} \frac{1}{n^{k}}\) . By (1) (VI, p.~283) one thus has, for \(|z| < 2\pi\)

\[
\frac {z}{e ^ {z} - 1} = 1 - \frac {z}{2} + \sum_ {n = 1} ^ {\infty} \frac {(- 1) ^ {n - 1} S _ {2 n}}{2 ^ {2 n - 1} \pi^ {2 n}} z ^ {2 n}\tag{13}
\]

whence we obtain the formula

\[
b _ {2 n} = (- 1) ^ {n - 1} (2 n)! \frac {2 \mathrm{S} _ {2 n}}{(2 \pi) ^ {2 n}} \quad \text {   for   } n \geqslant 1,\tag{14}
\]

a formula which shows in particular that the numbers \(S_{2n} / \pi^{2n}\) are rational. Clearly \(S_{k + 1} \leqslant S_k\) so, for every integer \(k \geqslant 2\) , we have \(S_k \leqslant S_2 = \pi^2 / 6 \leqslant 2\) ; from (14) one deduces the following inequalities for the Bernoulli numbers

\[
\frac {2 (2 n) !}{(2 \pi) ^ {2 n}} \leqslant | b _ {2 n} | \leqslant 4 \frac {(2 n) !}{(2 \pi) ^ {2 n}} \quad \text {   for   } n \geqslant 1.\tag{15}
\]

From these inequalities one can deduce an estimate for the Bernoulli polynomial \(\mathrm{B}_n(x) = \sum_{k=0}^{n}\binom{n}{k}b_kx^{n - k}\) ; in particular, for \(0 \leqslant x \leqslant 1\) one has

\[
\left| \mathrm{B} _ {n} (x) \right| \leqslant 4 \sum_ {k = 0} ^ {n} \binom {n} {k} \frac {k !}{(2 \pi) ^ {k}} = 4 \frac {n !}{(2 \pi) ^ {n}} \sum_ {k = 0} ^ {n} \frac {(2 \pi) ^ {k}}{k !} \leqslant 4 e ^ {2 \pi} \frac {n !}{(2 \pi) ^ {n}}.\tag{16}
\]

